\section{Conclusion and Future Work}
\label{sec:conclusion}

% Note: update the first sentence with the final headline result (whether
% the accuracy gain did or did not translate into a sequencing gain) once
% the experiments have been run.
We presented a reproducible programming-KT benchmark and a
knowledge-component / code-structure-aware KT variant, and studied not
only predictive accuracy but the extent to which accuracy translates into
better exercise sequencing---a coupling that is weak in general KT and
largely unexamined for programming. Our contributions span the pipeline
(RO1), the model (RO2), prediction evaluation (RO3), sequencing-utility
and correlation analysis (RO4), and cross-dataset/cross-domain
generalisation (RO5). The best-performing model and policy are the ones
adopted by the AlgoVenture adaptive engine (UC6, UC7).

\paragraph{Future work}
Directions include richer code representations, online evaluation with
real learners, and a live A/B study within AlgoVenture once the platform
is deployed to a pilot cohort.

\section*{Acknowledgment}
The authors thank their supervisor and the DSA teaching unit for guidance
and review.

\section*{AI Usage Declaration}
Generative AI tools were used to assist with drafting, editing, and code
scaffolding for this study. All AI-assisted content was reviewed and
verified by the authors. All models and datasets used are third-party,
publicly available artifacts (pyKT; CodeWorkout/CSEDM~2021; FalconCode;
ASSISTments/Junyi) and are cited accordingly. No experimental results
were produced by AI; all reported numbers come from running the
pipeline on the cited datasets.
